% !TeX root = ../../Book4.tex
% 纲要标题：### D.1 Frobenius 范畴与稳定范畴
% 纲要细目（写作范围）：
% #### D.1.1 正合结构中的投射与内射
% #### D.1.2 把经过投射内射对象的态射视为零
% #### D.1.3 合冲与余合冲函子

\section{Frobenius 范畴与稳定范畴}
\label{b4:sec:D01}

给一个模选取投射满射后，核记录生成元之间的关系；选取内射嵌入后，余核则记录另一方向的缺失。这两个对象通常依赖选择。如果把经过投射或内射对象的映射视为零，比较映射的差别就可能消失。Frobenius 范畴使这两种消去在同一个商范畴中进行。

\subsection{正合结构中的投射与内射}
\label{b4:subsec:D01:01}

\cref{b4:def:B-exact-category}中的正合范畴只指定一部分核余核对。因此，投射性与内射性也应当针对这些可容许短正合列定义。

\begin{definition}\label{b4:def:D-relative-projective-injective}
设 \(\mathcal E\) 为局部小正合范畴。若对象 \(P\in\mathcal E\) 对每个可容许满态射 \(p:B\to C\) 都使
\[
\operatorname{Hom}_{\mathcal E}(P,B)\longrightarrow
\operatorname{Hom}_{\mathcal E}(P,C)
\]
满射，就称 \(P\) 为该正合结构中的投射对象。若对象 \(I\) 对每个可容许单态射 \(i:A\to B\) 都使
\[
\operatorname{Hom}_{\mathcal E}(B,I)\longrightarrow
\operatorname{Hom}_{\mathcal E}(A,I)
\]
满射，就称 \(I\) 为该正合结构中的内射对象。若每个对象 \(X\) 都接收一个从投射对象出发的可容许满态射，称 \(\mathcal E\) 有足够多投射对象；若每个对象都经可容许单态射嵌入一个内射对象，称它有足够多内射对象。
\end{definition}

投射性说的是所给态射可以提升，内射性说的是所给态射可以延拓。它们的有限直和以及范畴内已有的直和项仍有相同性质，因为可以逐分量提升或延拓。在交换范畴的通常正合结构中，这些定义与模论中的定义相同；在仅含分裂列的正合结构中，每个对象都投射且内射。

\begin{definition}\label{b4:def:D-frobenius-category}
设 \(\mathcal E\) 为局部小正合范畴。如果它有足够多投射对象和足够多内射对象，而且两类对象相同，就称 \(\mathcal E\) 为 \term[frobenius-fanchou]{Frobenius 范畴}[Frobenius category]。把这类共同对象组成的满子范畴记为 \(\mathcal P\)，称其中的对象为投射内射对象。
\end{definition}

同一加性范畴上的不同正合结构可能给出不同的 \(\mathcal P\)。例如，对任意域 \(k\)，在双数代数 \(A=k[\varepsilon]/(\varepsilon^2)\) 的有限维左模范畴中，\(S=A/(\varepsilon)\) 在通常正合结构下不投射，在分裂正合结构下却投射且内射，具体比较见\cref{b4:prop:D-exact-structure-stable-comparison}。下面的商构造必须连同正合结构一起理解。Frobenius 范畴的定义与基本例子见 Bühler\cite[Section 13.4]{Buehler2010ExactCategories}。

\subsection{把经过投射内射对象的态射视为零}
\label{b4:subsec:D01:02}

\begin{proposition}\label{b4:prop:D-projective-ideal}
设 \(\mathcal E\) 为局部小 Frobenius 范畴，\(\mathcal P\) 为其投射内射对象组成的满子范畴。对 \(X,Y\in\mathcal E\)，令
\[
\mathcal P(X,Y)=
\{f:X\to Y:\ f\;\text{可以分解为}\;X\to P\to Y, P\in\mathcal P\}.
\]
则 \(\mathcal P(X,Y)\) 是 \(\operatorname{Hom}_{\mathcal E}(X,Y)\) 的子群，而且与任意态射在左右两侧复合后仍属于相应子群。
\end{proposition}
\symidx{\(\mathcal P(X,Y)\)：经过投射内射对象分解的态射子群}
\begin{proof}
零映射通过零对象分解，负映射通过同一个 \(P\) 分解。若 \(f=ba\)、\(f'=b'a'\) 分别通过 \(P,P'\)，则
\[
f+f'=
\begin{pmatrix}b&b'\end{pmatrix}
\begin{pmatrix}a\\a'\end{pmatrix}
\]
通过 \(P\oplus P'\) 分解，而有限直和仍投射且内射。对 \(u:X'\to X\)、\(v:Y\to Y'\)，有 \(vfu=(vb)(au)\)，仍通过 \(P\)。
\end{proof}

\begin{definition}\label{b4:def:D-stable-category}
设 \(\mathcal E\) 为局部小 Frobenius 范畴，\(\mathcal P\) 为其投射内射对象组成的满子范畴。保留 \(\mathcal E\) 的全部对象，规定
\[
\operatorname{Hom}_{\underline{\mathcal E}}(X,Y)
=\operatorname{Hom}_{\mathcal E}(X,Y)/\mathcal P(X,Y),
\qquad [g]\circ[f]=[g\circ f],
\]
所得范畴称为 \(\mathcal E\) 的\term[wendingfanchou]{稳定范畴}[stable category]，记为 \(\underline{\mathcal E}\)。映射 \(f\) 在此商中的类称为稳定态射。
\end{definition}
\symidx{\(\underline{\mathcal E}\)：Frobenius 范畴的稳定范畴}

\cref{b4:prop:D-projective-ideal}保证复合良定义。原来的零对象和有限双积仍给出商范畴的零对象和有限双积，所以 \(\underline{\mathcal E}\) 是局部小加性范畴。这里商掉的是态射子群，既没有把对象逐个删除，也没有反演某一类态射。

\begin{proposition}\label{b4:prop:D-stable-zero}
设 \(\mathcal E\) 为局部小 Frobenius 范畴。对象 \(X\) 在 \(\underline{\mathcal E}\) 中同构于零，当且仅当 \(X\) 是投射内射对象。特别地，若 \(P\) 投射且内射，则自然包含 \(X\to X\oplus P\) 在稳定范畴中可逆。
\end{proposition}
\begin{proof}
若 \(X\in\mathcal P\)，恒等映射通过 \(X\) 分解，故 \([1_X]=0\)。反过来，\([1_X]=0\) 给出 \(1_X=ba\)，其中 \(a:X\to P\)、\(b:P\to X\)、\(P\in\mathcal P\)。给定可容许满态射 \(q:E\to Y\) 及 \(f:X\to Y\)，先由 \(P\) 的投射性把 \(fb\) 提升为 \(\widetilde f:P\to E\)，满足 \(q\widetilde f=fb\)。于是
\[
 q(\widetilde f a)=fba=f,
\]
所以 \(\widetilde f a\) 提升 \(f\)，证明 \(X\) 投射。对可容许单态射 \(j:Y\to Z\) 及 \(g:Y\to X\)，由 \(P\) 的内射性把 \(ag\) 延拓为 \(\widetilde g:Z\to P\)，满足 \(\widetilde g j=ag\)。此时
\[
 (b\widetilde g)j=bag=g,
\]
证明 \(X\) 内射。这里已经有实际的收缩数据 \(a,b\)，不要求 \(P\) 的互补对象存在，也不要求范畴幂等完备。最后，包含与投影的一个复合是恒等，另一个与恒等之差只作用于 \(P\) 分量，故通过 \(P\) 分解。
\end{proof}

\subsection{合冲与余合冲函子}
\label{b4:subsec:D01:03}

在 Frobenius 范畴中，为每个对象固定可容许短正合列
\begin{equation}\label{b4:eq:D-syzygy-sequences}
0\longrightarrow\Omega X\xrightarrow{j_X}P_X\xrightarrow{p_X}X\longrightarrow0,
\qquad
0\longrightarrow X\xrightarrow{i_X}I_X\xrightarrow{q_X}\Sigma X\longrightarrow0,
\end{equation}
其中 \(P_X,I_X\in\mathcal P\)。称 \(\Omega X\) 为这次选择的合冲对象，\(\Sigma X\) 为余合冲对象。对大范畴，这些逐对象指定作为已选数据，沿用\subsecref{b4:subsec:0101:03}的大小与选择约定。
\symidx{\(\Omega X,\ \Sigma X\)：稳定范畴中的合冲与余合冲对象}

\begin{theorem}\label{b4:thm:D-stable-syzygy-equivalence}
设 \(\mathcal E\) 为局部小 Frobenius 范畴，并固定式\eqref{b4:eq:D-syzygy-sequences}中的序列。核与余核的比较映射分别定义加性函子
\[
\Omega,\Sigma:\underline{\mathcal E}\longrightarrow\underline{\mathcal E}.
\]
更换已选序列得到的函子与原函子自然同构，而且
\[
\Omega\Sigma\cong\operatorname{Id}_{\underline{\mathcal E}},
\qquad
\Sigma\Omega\cong\operatorname{Id}_{\underline{\mathcal E}}.
\]
因此 \(\Sigma\) 是自等价，\(\Omega\) 是其拟逆。
\end{theorem}
\begin{proof}
给定 \(f:X\to Y\)，投射性提供 \(\widetilde f:P_X\to P_Y\)，满足 \(p_Y\widetilde f=fp_X\)。固定这张提升后，核的泛性质才给出唯一 \(a:\Omega X\to\Omega Y\)，使 \(j_Ya=\widetilde f j_X\)；提升本身通常不唯一。若换成另一提升 \(\widetilde f'\)，两提升之差唯一写成 \(j_Yc\)，其中 \(c:P_X\to\Omega Y\)。由 \(j_Y(a-a')=j_Ycj_X\) 及 \(j_Y\) 为单态射，得到 \(a-a'=cj_X\)。这个差通过 \(P_X\) 分解，故 \([a]=[a']\)。构造依赖提升的选择，但商去这项误差后，稳定类不再依赖它。

若 \(f\) 通过投射对象 \(Q\) 分解为 \(X\xrightarrow{u}Q\xrightarrow{v}Y\)，先把 \(v\) 提升为 \(\widetilde v:Q\to P_Y\)。取 \(\widetilde f=\widetilde vup_X\)，则 \(\widetilde f j_X=0\)，故所诱导的核映射为零。这说明构造只依赖 \([f]\)。两提升的复合是复合态射的一种提升，提升之和是态射之和的一种提升，恒等可以用恒等提升；由刚证的独立性，得到加性函子 \(\Omega\)。

内射性同样提供 \(\widehat f:I_X\to I_Y\)，满足 \(\widehat f i_X=i_Yf\)。余核的泛性质给出 \(b:\Sigma X\to\Sigma Y\)。两个延拓之差写成 \(c q_X\)，其中 \(c:\Sigma X\to I_Y\)，所以两个余核映射之差为 \(q_Yc\)，通过 \(I_Y\) 分解。若 \(f\) 通过内射对象 \(Q\) 分解，先把 \(X\to Q\) 沿 \(i_X\) 延拓，再接 \(Q\to Y\to I_Y\)，所得余核映射为零。复合、加法与恒等的相容性由同样的比较得到，因此 \(\Sigma\) 也是加性函子。

现在比较同一对象 \(X\) 的两套投射序列，记其满态射为 \(p_X:P_X\to X\) 与 \(p'_X:P'_X\to X\)。投射性给出 \(u:P_X\to P'_X\)、\(v:P'_X\to P_X\)，满足 \(p'_Xu=p_X\)、\(p_Xv=p'_X\)，从而诱导两套核之间的比较映射。其复合由 \(vu\) 诱导，而 \(p_Xvu=p_X\)，所以 \(vu\) 与 \(1_{P_X}\) 都是 \(1_X\) 的提升。前面已证明两个提升所诱导的核映射之差通过 \(P_X\) 分解；恒等提升诱导恒等核映射，因此该复合在稳定范畴中为恒等，反向同理。对 \(f:X\to Y\)，两条比较路径都来自 \(f\) 的提升；其差通过相应投射对象分解，故这些比较同构自然。内射序列的比较由上述延拓论证得到相同结论。

第二条序列原来是 \(X\) 的内射序列；因 \(I_X\) 同时投射，它又可作 \(\Sigma X\) 的投射序列，其核就是 \(X\)。构造 \(\Sigma[f]\) 时的延拓 \(\widehat f:I_X\to I_Y\) 满足 \(q_Y\widehat f=(\Sigma f)q_X\)，所以也可作 \(\Sigma f\) 的投射提升。它又满足 \(\widehat f i_X=i_Yf\)，而 \(i_Y\) 为单态射，故限制在核 \(X\) 上所诱导的映射恰为 \(f\)。因此用这组序列计算 \(\Omega\Sigma\) 得到恒等函子，再与固定的投射序列比较，就得到自然同构 \(\Omega\Sigma\cong\operatorname{Id}\)。

第一条序列原来是 \(X\) 的投射序列；因 \(P_X\) 同时内射，它又可作 \(\Omega X\) 的内射序列，其余核就是 \(X\)。对 \(f\) 使用构造 \(\Omega[f]\) 时的提升图，余核上的映射也恰为 \(f\)。用这组序列计算 \(\Sigma\Omega\) 因而得到恒等函子，再与固定的内射序列比较，得到自然同构 \(\Sigma\Omega\cong\operatorname{Id}\)。两种同构的自然性都来自刚才的序列比较，而不只是在每个对象上单独得到同构。
\end{proof}

核和余核在原范畴中可以不同；自然同构发生在稳定范畴中。投射内射项同时能消去两种比较误差，正是合冲与余合冲成为互逆操作的原因。

\begin{proposition}\label{b4:prop:D-exact-structure-stable-comparison}
设 \(k\) 为域，\(A=k[\varepsilon]/(\varepsilon^2)\)，\(S=A/(\varepsilon)\)，并在有限维幺性左 \(A\)-模的同一个加性范畴上分别指定通常短正合结构和分裂短正合结构。两者都是 Frobenius 正合结构，但所得稳定范畴不同。

在通常结构中，投射内射对象恰为有限自由模 \(A^r\)，且
\[
 0\longrightarrow S\xrightarrow{\iota}A\xrightarrow{\pi}S
 \longrightarrow0,\qquad\iota(\overline1)=\varepsilon,
\]
是不分裂的短正合列，其中 \(\pi\) 为商映射。此时
\[
 \operatorname{End}_{\underline{A\text{-}\mathrm{mod}}}(S)\cong k,
\]
所以 \(S\) 在稳定范畴中不为零。在分裂结构中，每个对象都投射且内射，稳定范畴等价于零范畴。
\end{proposition}
\begin{proof}
\subsecref{b4:subsec:D02:02}用取 \(\varepsilon\) 系数的泛函证明左模 \(A\cong D(A)\)，由\cref{b4:prop:D-frobenius-module-category}，通常结构为 Frobenius 结构。\cref{b4:prop:D-dual-number-stable}给出全部有限维模的分解 \(A^r\oplus S^s\)，并证明投射内射对象恰为 \(A^r\)。该证明还说明 \(\pi:A\to S\) 没有 \(A\)-线性截面，所以所列短正合列不分裂。任何 \(S\to A^r\) 的像落在 \(\varepsilon A^r\)，任何 \(A^r\to S\) 又杀掉它，因而经过投射内射对象的 \(S\to S\) 映射全为零。因此稳定自同态环仍是 \(\operatorname{End}_A(S)\cong k\)，其恒等类非零。

只含分裂列的结构由\cref{b4:prop:B-split-exact}给出。对分裂满态射 \(q:B\to C\)，取截面 \(s:C\to B\)，则任意 \(f:X\to C\) 都由 \(sf\) 提升。对分裂单态射 \(j:B\to C\)，取收缩 \(r:C\to B\)，则任意 \(g:B\to X\) 都由 \(gr\) 延拓。所以每个对象都投射且内射，恒等态射本身提供足够多投射与内射对象。每个态射 \(f:X\to Y\) 都通过投射内射对象 \(X\) 分解，稳定 Hom 群全部为零，所得范畴遂等价于零范畴。
\end{proof}
